\section{Semi-analytic curved-tube reference (SACR)}
\label{sec:sacr}
SACR evaluates the free-space electrostatic field of a prescribed curved charge
distribution without a three-dimensional mesh. The target geometry is a 5 m
beam on a 2 m bend radius, spanning $143.24^\circ$, with 5 mm radial extent.
This section uses an ellipsoidally tapered source to obtain an exact Dubin
limit. \Cref{sec:frozen} extends the field comparison to a constant-section
tube; neither test establishes its self-consistent tracking accuracy.

\begin{figure}[H]
 \centering
 % Prescribed SACR sources and their distinct straight limits; no solver grid.
% The left envelope is u=+/-a0*sqrt(1-s^2/b^2), mapped onto a circular path.
% The midpoint source patch has a'=c'=0, so its outward normal is radial.
% Transverse dimensions are deliberately enlarged. Connections are geometry,
% not computed electric-field vectors or magnitudes.
\begin{tikzpicture}[
 x=1cm,y=1cm,font=\small,>=Latex,
 srsource/.style={blue!65!black,line width=.8pt},
 srthin/.style={black!45,line width=.45pt},
 sraxis/.style={->,black!70,line width=.6pt},
 srdim/.style={<->,black!65,line width=.55pt},
 srlabel/.style={fill=white,inner sep=2pt},
 srtarget/.style={fill=orange!85!black,draw=white,line width=.6pt}
]
 \path[use as bounding box] (0,0) rectangle (16.6,10.8);
 \node[anchor=north west,font=\bfseries\small] at (.1,10.65)
   {(a) Mesh-free curved-source integration};
 \node[anchor=north west,black!65] at (.1,10.1)
   {Prescribed geometry; uniform physical $\rho$};
 \node[anchor=north west,black!65] at (.1,9.65)
   {Orbit-plane view; $y$ out of the page};
 \node[anchor=north west,font=\bfseries\small] at (8.7,10.65)
   {(b) Two source families, two straight limits};

 % Circular reference: 5 m source / 2 m radius = 143.239449 degrees.
 \coordinate (srO) at (1.35,6.0);
 \fill[blue!9]
  plot[domain=-71.619724:71.619724,samples=121]
   ({1.35+(2.65-.22*sqrt(max(0,1-(\x/71.619724)^2)))*cos(\x)},
    {6.0+(2.65-.22*sqrt(max(0,1-(\x/71.619724)^2)))*sin(\x)})
  -- plot[domain=71.619724:-71.619724,samples=121]
   ({1.35+(2.65+.22*sqrt(max(0,1-(\x/71.619724)^2)))*cos(\x)},
    {6.0+(2.65+.22*sqrt(max(0,1-(\x/71.619724)^2)))*sin(\x)}) -- cycle;
 \foreach \side in {-1,1} {
  \draw[srsource]
   plot[domain=-71.619724:71.619724,samples=121]
    ({1.35+(2.65+\side*.22*sqrt(max(0,1-(\x/71.619724)^2)))*cos(\x)},
     {6.0+(2.65+\side*.22*sqrt(max(0,1-(\x/71.619724)^2)))*sin(\x)});
 }
 \draw[srthin,dashed] ($(srO)+(-71.619724:2.65)$)
  arc[start angle=-71.619724,end angle=71.619724,radius=2.65];
 \foreach \theta in {-60,-45,-25,25,45,60} {
  \pgfmathsetmacro{\srwidth}{.22*sqrt(1-(\theta/71.619724)^2)}
  \fill[blue!65!black] ($(srO)+(\theta:{2.65+\srwidth})$) circle[radius=1.2pt];
  \fill[blue!65!black] ($(srO)+(\theta:{2.65-\srwidth})$) circle[radius=1.2pt];
 }

 % Source patch at the maximum cross-section: its normal is exactly radial.
 \coordinate (srP) at (4.22,6.0);
 \draw[blue!70!black,line width=2.5pt]
  plot[domain=-4:4,samples=15]
   ({1.35+(2.65+.22*sqrt(1-(\x/71.619724)^2))*cos(\x)},
    {6.0+(2.65+.22*sqrt(1-(\x/71.619724)^2))*sin(\x)});
 \fill[blue!65!black] (srP) circle[radius=1.7pt];
 \draw[->,blue!65!black,line width=.8pt] (srP) -- (5.55,6.0)
  node[right,black] {$\bm n'\,\dd S'$};
 \node[anchor=north west,srlabel,align=left] at (4.35,5.7)
  {Source patch $\bm X'$\\$s'=0$};
 \coordinate (srT) at ($(srO)+(25:2.65)$);
 \draw[black!65,line width=.6pt] (srP) -- (srT)
  node[midway,left=3pt,srlabel] {$d$};
 \fill[srtarget] (srT) circle[radius=2.7pt];
 \node[anchor=west,srlabel] at ($(srT)+(.22,.14)$)
  {Target $\bm X$};
 \draw[srdim] ($(srO)+(.08,0)$) -- (3.85,6.0)
  node[midway,above=3pt,srlabel] {$R=2\,\mathrm m$};
 \fill[black!65] (srO) circle[radius=1.1pt];
 \node[anchor=north,align=center] at ($(srO)+(0,-.16)$)
  {Curvature\\centre};
 \node[anchor=south,srlabel] at ($(srO)+(71.619724:2.65)+(0,.12)$)
  {$s=+b$};
 \node[anchor=north,srlabel] at ($(srO)+(-71.619724:2.65)+(0,-.08)$)
  {$s=-b$};

 % A local radial/tangential basis away from the integration target.
 \coordinate (srB) at ($(srO)+(-42:2.65)$);
 \draw[sraxis] (srB) -- ($(srB)+(-42:.65)$) node[below right] {$u$};
 \draw[sraxis] (srB) -- ($(srB)+(48:.65)$)
  node[above left,srlabel,inner sep=.5pt] {$s$};
 \node[anchor=west,align=left] at (.15,2.65)
  {$L=2b=5\,\mathrm m$\\$L/R=143.24^\circ$};
 \node[anchor=west,align=left,black!65] at (.15,1.45)
  {Surface quadrature points are schematic\\
   Transverse thickness enlarged; no volume mesh};
 \node[anchor=west] at (.15,.45)
  {$d=|\bm X-\bm X'|$: physical separation};

 % Family 1: fixed material semiaxes, curvature tends to zero.
 \node[align=center] at (12.7,9.45)
  {Ellipsoidally tapered cross-section};
 \draw[sraxis] (8.9,8.0) -- (9.95,8.0)
  node[midway,above=3pt,srlabel] {$R\to\infty$};
 \draw[srsource,fill=blue!9] (12.7,8.0) ellipse[x radius=2.25,y radius=.57];
 \draw[srthin,dashed] (12.7,8.0) ellipse[x radius=2.25,y radius=.17];
 \draw[sraxis] (10.2,8.0) -- (15.45,8.0) node[right] {$s$};
 \node[anchor=north] at (10.45,7.25) {$-b$};
 \node[anchor=north] at (14.95,7.25) {$+b$};
 \node[align=center] at (12.7,6.4)
  {$a(s),c(s)\ \propto\ \sqrt{1-s^2/b^2}$\\[3pt]
   Dubin ellipsoid: exact straight limit};

 % Family 2: the same limit does not turn a constant-section source into an ellipsoid.
 \node[align=center] at (12.7,4.85)
  {Constant cross-section; sharp ends};
 \draw[sraxis] (8.9,3.4) -- (9.95,3.4)
  node[midway,above=3pt,srlabel] {$R\to\infty$};
 \fill[blue!9] (10.45,2.85) rectangle (14.95,3.95);
 \draw[srsource] (10.45,2.85) -- (14.95,2.85);
 \draw[srsource] (10.45,3.95) -- (14.95,3.95);
 \draw[srsource] (10.45,3.4) ellipse[x radius=.16,y radius=.55];
 \draw[srsource,fill=blue!9] (14.95,3.4) ellipse[x radius=.16,y radius=.55];
 \draw[srthin,dashed] (10.45,3.4) -- (14.95,3.4);
 \node[align=center] at (12.7,2.1)
  {Finite elliptic cylinder\\[3pt]Not the Dubin ellipsoid};
 \node[align=center,black!65,text width=7.5cm] at (12.7,.7)
  {Straight limits hold at fixed source lengths,\\
   cross-section parameters and charge};
\end{tikzpicture}

 \caption{SACR is a mesh-free Coulomb reference for a prescribed physical
 charge distribution. Source-surface patches contribute through their
 outward vector areas and physical source-target separations. At fixed
 semiaxes and charge, $R\to\infty$ maps the ellipsoidally tapered family to
 Dubin's straight ellipsoid; a constant-section tube instead tends to a
 finite elliptic cylinder. Surface-only integration shown here assumes
 uniform physical density; a varying longitudinal density also requires
 the volume-gradient term. Transverse dimensions are enlarged.}
 \label{fig:sacr-reference-geometry}
\end{figure}
The finite-radius source is retained in the Coulomb distance. Its geometry
is integrated directly, without the IPPL grid or the first-order transverse
curvature expansion used by the OPALX OPEN backend. The two straight limits
in the illustration therefore describe different physical reference sources.
The reference-path coordinate map is orthogonal but not conformal: its scale
factors are $1,1,1+u/R$. The curved fields are evaluated from physical Coulomb
distances, rather than obtained by a conformal transformation of a straight
beam solution.
This describes the present method; it does not rule out other mapped
formulations, which must still represent the physical three-dimensional
source and its longitudinal coupling.

\clearpage
\subsection{Field model and straight limit}
For reference arclength $s$, outward radial displacement $u$ and vertical
coordinate $y$, the OPALX laboratory position and volume element are
\begin{align}
 \bm X(s,u,y)&=\big((R+u)\cos(s/R)-R,\ y,\ (R+u)\sin(s/R)\big)^\T,\\
 \dd V&=\left(1+\frac{u}{R}\right)\dd s\,\dd u\,\dd y.
 \label{eq:sacr-map}
\end{align}
For uniform physical volume density $\rho$, the Coulomb field has the exact
surface form
\begin{equation}
 \bm E(\bm x)=\frac{\rho}{4\pi\eps}
 \int_V\frac{\bm x-\bm x'}{|\bm x-\bm x'|^3}\,\dd V'
 =\frac{\rho}{4\pi\eps}
 \oint_{\partial V}\frac{\bm n'(\bm x')}{|\bm x-\bm x'|}\,\dd S'.
 \label{eq:sacr-surface}
\end{equation}
The second equality follows from the divergence theorem, with outward source
normal $\bm n'$. Two-dimensional boundary quadrature gives all three field
components, without a fitted $g$ or softening length.

For the Dubin comparison, choose
\begin{equation}
 a(s)=a_0\sqrt{1-s^2/b^2},\qquad
 c(s)=c_0\sqrt{1-s^2/b^2},\qquad |s|\le b,\qquad L=2b.
 \label{eq:sacr-taper}
\end{equation}
The physical density is constant, while the line density is parabolic.
The volume remains $4\pi a_0bc_0/3$ under bending: the term odd in $u$ in
\cref{eq:sacr-map} integrates to zero. At fixed semiaxes and charge,
$R\to\infty$ gives the Cartesian ellipsoid and its exact field
$E_i=\rho D_i x_i/\eps$ from Dubin~\cite{dubin}. Our $D_i$ equal half of
Dubin's coefficients, so $\sum_iD_i=1$.

A constant-section tube instead approaches a finite cylinder as $R\to\infty$;
it is not the finite ellipsoid used in this check. The 0.5 mm vertical semiaxis
in the thin verification cases is a benchmark choice, not a newly specified
real-beam parameter.

For an independent tangential check, let $A$ be the transverse ellipse
$u^2/a_0^2+y^2/c_0^2\le1$. At target $(u_t,y_t,s_t)$ define
$r_t=R+u_t$, $r_s=R+u$ and
$s_\pm=\pm b\sqrt{1-u^2/a_0^2-y^2/c_0^2}$. Exact integration along each
source filament gives
\begin{equation}
 E_s=\frac{\rho}{4\pi\eps}\int_A\frac{r_s}{r_t}
 \left(\frac{1}{d_+}-\frac{1}{d_-}\right)\dd u\,\dd y,
 \qquad d_\pm=|\bm X(s_t,u_t,y_t)-\bm X(s_\pm,u,y)|.
 \label{eq:sacr-endpoint}
\end{equation}
The straight limit has $r_s/r_t\to1$; the tapered source requires these
transverse-dependent ends.

\subsection{Error measures}
Fields are compared at matched material coordinates $\bm\xi=(u,y,s)$ and projected
onto each target's local radial, vertical and tangential basis. Let
$E_i^{\rm SACR}(R;\bm\xi)$ denote that local field and
$\ell_i=(a_0,c_0,b)_i$. Define
\begin{equation}
 E_{i,*}=\frac{\rho}{\eps}D_i\ell_i,\qquad
 \delta_i(R,\bm\xi)=\frac{E_i^{\rm SACR}(R;\bm\xi)-E_i^{\rm Dubin}(\bm\xi)}{E_{i,*}}.
 \label{eq:sacr-error}
\end{equation}
This component scale retains targets where the exact field vanishes. Changes
under quadrature refinement measure numerical resolution; finite-$R$
differences from Dubin also contain the physical curvature effect.

\clearpage
\subsection{Multi-point straight-limit checks}
The study uses 33 interior points per geometry: the centre, axis and midplane
points, signed off-axis points, and points near the ends and boundary. The
largest normalized ellipsoidal radius is 0.911. Three geometries are checked:
a 5 mm sphere, the previous $(5,50,0.5)$ mm ellipsoid, and a
$(5,2500,0.5)$ mm ellipsoid. Each value in the table is a maximum over all
33 points, expressed as a percentage of the corresponding component scale.
\tabfile{sacrflat}
\reportfigure{sacr-flat-limit}{Straight-limit SACR fields versus Dubin at all
33 points per geometry. The dashed identity line is the exact ellipsoid
field. Lower panels retain signed residuals, including zero-field targets;
normalization uses $E_{i,*}$ rather than the local field value.}{fig:sacr-flat}
Angular quadrature starts at 128 nodes and refines to at most 4096 where
needed. The largest measured straight-limit difference is
$3.97\times10^{-14}E_{i,*}$. The largest last-refinement change over the full
study is $3.64\times10^{-9}E_{i,*}$, below the requested $10^{-8}E_{i,*}$
criterion. These are sampled field checks, not a uniform error bound over the
entire charge distribution.
Dubin's coefficients set component scales and quadrature stopping criteria;
the fields themselves come from the Coulomb integral. A separate scalar
implementation agrees with the vectorized calculation within
$1.68\times10^{-14}E_{i,*}$.

\clearpage
\subsection{Radius convergence to Dubin}
For the 5 m tapered source, $a_0=5$ mm, $b=2.5$ m and $c_0=0.5$ mm are fixed
while $R$ varies from 2 to 2000 m and infinity. The same 33 material points
are used at each radius. The finite-radius difference contains the physical
change from a curved source to a straight ellipsoid; it is not an SACR error.
\reportfigure{sacr-radius-limit}{Convergence of the 5 m SACR source toward the
straight Dubin field. Solid curves show $\max_{\bm\xi}|\delta_i|$ over the 33
points; dashed curves show the separate last quadrature-refinement changes.
The $R=\infty$ checks are reported in the preceding table, not placed at an
artificial finite coordinate on this logarithmic axis.}{fig:sacr-radius}
\tabfile{sacrradius}
At $R=2$ m, the largest curvature differences are 0.3859\% radial,
0.4112\% tangential and 0.002801\% vertical, on the component scales defined
in \cref{eq:sacr-error}. All three decrease toward the straight limit. At
2000 m the vertical difference remains about 7.7 times the last-refinement
change; no universal $1/R^2$ rate is assumed.

An independent tangential endpoint integral agrees with the surface field
within $6.97\times10^{-15}E_{s,*}$ across 99 comparisons. Symmetry checks retain
the expected tangential/vertical zeros. At three interior points with $R=2$ m,
Cartesian finite differences use $\bm F=4\pi\eps\bm E/\rho$, whose exact
interior identities are $\nabla\cdot\bm F=4\pi$ and
$\nabla\times\bm F=\bm0$. The maximum errors are
$|\nabla\cdot\bm F/(4\pi)-1|=5.20\times10^{-12}$ and
$|\nabla\times\bm F|/(4\pi)=5.23\times10^{-12}$.
An independent volume integral at the bend centre, outside the source, agrees
with the boundary integral within $6.88\times10^{-13}$ relative.

This validates SACR's static field operator on the tested tapered family.
\Cref{sec:frozen} treats constant-section fields and their OPALX comparison.
A fixed prescribed tube does not supply a tracking reference after space
charge deforms its density.

